% Осень26!B32; нумерация внутри столбца без пустых строк.
% Источник: лекция 6, §4 (строка 198)
\definitionticket{29}{Взвешенный образ меры}
\label{definition:29}

\begin{definition}
Пусть $(X,\mathcal A,\mu)$ --- пространство с мерой, $(Y,\mathcal B)$ --- множество с $\sigma$-алгеброй, $\varphi:X\to Y$ --- измеримое отображение, а $\omega:X\to[0,+\infty]$ --- измеримая функция.
\textbf{Взвешенным образом меры $\mu$ с весом $\omega$ при отображении $\varphi$} называется мера $\nu$ на $\mathcal B$, заданная формулой
\[
  \nu(E):=\int_{\varphi^{-1}(E)}\omega\,d\mu,
  \qquad E\in\mathcal B.
\]
\end{definition}
